\subsubsection{SEM Description and Theory}\label{SEC:SEM:theory}

For turbulent flow, one rarely knows what precise unsteady values to
specify at the inflow boundary.
%
Often, however, time-averaged boundary-layer profile data exists or can be
estimated based on Reynolds number.
%
One way to provide an unsteady turbulent inflow is to add
an upstream region of grid which has a steady uniform inflow,
some numerical or physical tripping mechanism, and
a long run of length for the boundary layer to grow to the desired
thickness before entering the region of interest.
%
Unfortunately, this method is computationally expensive due to the added
upstream length and unsteady iterations needed to obtain time-averaged
statistics.
%
It also involves numerous parameters related to the trip and run length
that likely require adjustment and repetition of the whole simulation.
%
Furthermore, the method might not even provide an adequate boundary layer
due to other effects such as upstream surface curvature or pressure gradient.

%%%%%%%%%%%%%%%%%%%%%%%%%%%%%%%%%%%%%%%%%%%%%%%%%%%%%%%%%%%%%%%%%%%%%%%%%%%%%%%%

The synthetic eddy method (SEM)\cite{Jarrin_2006} allows the user to specify a
time-averaged inflow profile for the mean flow and turbulent stresses.
%
The algorithm will then use that data to artificially generate unsteady inflow
values that, when averaged over a sufficient amount of time, will replicate
the specified profiles.
%
The way in which the SEM works can be summarized as follows.
%
\begin{enumerate}
\item
The SEM profile file is read to obtain the mean flow and turbulent statistics
to be imposed on the inflow boundary.
%
The height of the SEM profile is determined from the furthest point
specified in the file.
%
\item
A grid-free virtual upstream volume is created that overlaps the
physical inflow plane in the streamwise direction.
%
\item
Synthetic eddies with random intensities are randomly seeded within
this virtual volume.
%
\item
Each grid point on the physical inflow boundary computes its own unique
fluctuating velocity based on a weighted aggregate of the intensities
of the synthetic eddies within the specified turbulent length scale.
%
\item
The fluctuating velocity is then combined with the user-provided
time-averaged profile for the mean flow and turbulent stresses to
create an unsteady turbulent inflow boundary condition.
%
\item
The eddies within the virtual volume are convected as particles
using the specified convective velocity in the direction perpendicular to the
physical inflow plane and updated each time step.
%
Every time an eddy leaves the virtual volume, a new eddy,
with a random intensity, is created at a random location on the inflow face
of the virtual volume.
%
\item
Repeat the process, starting from step 4.
\end{enumerate}

%%%%%%%%%%%%%%%%%%%%%%%%%%%%%%%%%%%%%%%%%%%%%%%%%%%%%%%%%%%%%%%%%%%%%%%%%%%%%%%%

The following sections, \ref{SEC:SEM:theory:Nomenclature} --
\ref{SEC:SEM:theory:ConvectEddies}, provide additional details on
the SEM implementation in GFR.

%###############################################################################

\subsubsubsection{SEM Nomenclature}
\label{SEC:SEM:theory:Nomenclature}

%%%%%%%%%%%%%%%%%%%%%%%%%%%%%%%%%%%%%%%%%%%%%%%%%%%%%%%%%%%%%%%%%%%%%%%%%%%%%%%%

In this section, we define a number of variables that will be used
quite often when describing the SEM.
%
Since multiple SEM inflow boundaries can be specified in the
\inp[nmlBC]{nml\_bc} namelist input block, we will use $i$ to represent
a specific \BCinptext[i]{SEM} instance.\\[12pt]
%
\begin{tabular}{%
      @{\hspace{2.0em}} l%
      @{\hspace{1.0ex}} c%
      @{\quad} >{\raggedright\arraybackslash}p{0.8\linewidth}%
      @{}}
  %
  \uBC & $\coloneq$ & The mean velocity at the freestream or edge
    of the boundary layer in the provided SEM profile file,
    \BCinpBC[nmlBC::BCInput/SEM/ProfileFile][i]{SEM\%Profile\_File}.
    See \sref{SEC:SEM:usage:profile} for details of the profile file.\\
  %
   & & \\
  %
  \lBC & $\coloneq$ & The height of the boundary layer within the SEM
    profile. It is important to note that \lBC{} is not necessarily the same
    as \hSEM{} defined below. Ideally, the SEM profile should include the
    entirety of the boundary layer, so in general, $\hSEM \geq \lBC$.
    While \lBC{} is used in our description of the boundary layer to be
    imposed, the SEM algorithm only uses \hSEM{}.\\
  %
   & & \\
  %
  \uSEM & $\coloneq$ & The SEM convective velocity, which is the speed
    at which the synthetic eddies are propagated through the SEM box. It is
    specified using the input
    \BCinpBC[nmlBC::BCInput/SEM/ConvectiveVelocity]%
            [i]{SEM\%Convective\_Velocity}.\\
  %
   & & \\
  %
  \lSEM & $\coloneq$ & The SEM length scale, which determines the
  \quoterm{zone of influence} or \quoterm{sensory distance} of each synthetic
  eddy. It is specified using the input
  \BCinpBC[nmlBC::BCInput/SEM/LengthScale][i]{SEM\%Length\_Scale}.\\
  %
   & & \\
  %
  \hSEM & $\coloneq$ & The height of the SEM profile. This is obtained
    from the furthest profile point specified in the SEM profile file.\\
  %
   & & \\
  %
  \key*[key:SEMVolume]{\vSEM}[Synthetic Eddy Method (SEM)!volume] & $\coloneq$ &
  The volume of the SEM box, defined by \eref{EQN:VolSEM}.
  \\
  %
   & & \\
  %
  \key*[key:EddyVolume]{\vEddy}[eddy volume] & $\coloneq$ &
  The volume of an eddy,
  $\lSEM^3$ for a cube or
  $\sfrac{\pi}{6} \; \lSEM^3$ for a sphere,
  depending on 
  \inp[nmlSEM::UseSphericalEddyVolume]{Use\_Spherical\_Eddy\_Volume}.
  %
\end{tabular}

%###############################################################################

\subsubsubsection{Create a virtual box at SEM inflow}
\label{SEC:SEM:theory:VirtualBox}

To summarize how the SEM works, we start with \fref{FIG:SEM:FlatPlate} that
gives a conceptual illustration of a flat plate case.
%
\begin{figure}[tb]
  \centering
  \insertimage[width=0.95\textwidth]{sem/traditional_inflow}
% \caption[Idealized flat plate inflow profile.]{To illustrate the concept of
% the SEM method, we are using a conceptual drawing of a simple flat plate
% case. The no-slip wall is the gray surface extending downstream. The blue
% plane is the inflow boundary surface, and the blue arrows show an idealized
% flow profile over this flat plate including the boundary layer near the
% wall.}
  \caption{Illustration of an inflow profile for a flat plate.}
  \label{FIG:SEM:FlatPlate}
\end{figure}
%
The no-slip wall is the bottom gray surface extending downstream,
and the blue plane on the left is the inflow boundary surface.
%
The blue arrows show an idealized flow profile over this flat plate including
the boundary layer near the wall.

%%%%%%%%%%%%%%%%%%%%%%%%%%%%%%%%%%%%%%%%%%%%%%%%%%%%%%%%%%%%%%%%%%%%%%%%%%%%%%%%

The SEM starts by creating a grid-free virtual volume\slash box that overlaps
the physical inflow plane in the streamwise direction.
%
In \Fref{FIG:SEM:box}, the green rectangular box overlapping the blue inflow
plane represents the virtual volume that is created for the SEM.
%
\begin{figure}[tb]
  \centering
  \insertimage[width=0.95\textwidth]{sem/sem_box}
% \caption[SEM box added to flat plate inflow plane.]{The green rectangular box
% overlapping the blue inflow plane represents the virtual volume that is
% created to house the SEM.}
  \caption{Illustration of the SEM virtual volume on flat plate inflow plane.}
  \label{FIG:SEM:box}
\end{figure}
%
In the streamwise direction ($x$ in this example), assuming the inflow plane is
located at a constant $x$ value of $x_\subBC$, the SEM box is axially
centered at $x_\subBC$ and extends upstream and downstream of the inflow
plane by a distance of the SEM length scale, \lSEM.
%
In the spanwise direction ($z$ in this example), the SEM box extends the full
span of the inflow plane (say from $z_{\textrm{min}}$ to $z_{\textrm{max}}$),
with an extra \lSEM{} added to both ends.
%
In the wall normal direction ($y$ in this example), the SEM box extends from the
wall surface, $y_{\textrm{wall}}$, to the height of the SEM profile, \hSEM, also
with an extra \lSEM{} added to both ends.
%
The final SEM box is the rectangular volume contained by 
%
\begin{gather}
  \begin{aligned}
    x_\subBC          - \lSEM & \leq x \leq x_\subBC          + \lSEM \\
    y_{\textrm{wall}} - \lSEM & \leq y \leq y_{\textrm{wall}} + \hSEM
                                                              + \lSEM \\
    z_{\textrm{min}}  - \lSEM & \leq z \leq z_{\textrm{max}}  + \lSEM
  \end{aligned}
\end{gather}
%
with volume of
\begin{gather}
  \label{EQN:VolSEM}
  \vSEM = (2\lSEM)(\hSEM+2\lSEM)(2\lSEM)
\end{gather}

%###############################################################################

\subsubsubsection{Define flow profile for SEM}
\label{SEC:SEM:theory:DefineProfile}

The SEM requires a data file containing mean flow and turbulent stress
variables for the boundary layer profile that the user wants to impose at
the SEM inflow.
%
Details regarding the required content and format of this data file are given in
\sref{SEC:SEM:usage:profile}.
%
Continuing with the previous flat plate example, \fref{FIG:SEM:InflowProfile}
shows the boundary layer section of the idealized flow profile, identified by
the green arrows, which will be imposed at the inflow.
%
\begin{figure}[tb]
  \centering
  \insertimage[width=0.95\textwidth]{sem/sem_profile}
  \caption[Illustration of the boundary layer section of the inflow profile.]%
  {Illustration of the boundary layer section of the inflow profile,
  shown in green.
  The mean flow and turbulent stress variables defining this boundary layer
  are given in a data file that is input to the SEM boundary condition.
  The SEM will use this profile information to create an unsteady boundary
  condition with velocity fluctuations that should recreate this profile when
  averaged over time.}
  \label{FIG:SEM:InflowProfile}
\end{figure}
%
This flow profile will contain the mean flow variables
%
\begin{equation*}
  \overline{\bm{q}} =
  \left[\overline{\rho},\overline{u}_i,\overline{T}\right]
\end{equation*}
%
at a set of locations within the boundary layer with an accompanying wall
distance variable.
%
The specific Reynolds stress tensor,
$R_{ij}=\ReyAve{\ReyFlx{u}_i\ReyFlx{u}_j}$,
also needs to be provided in the data file
to transform randomly generated numbers into velocity fluctuations that are
physically correlated to the turbulent boundary layer profile.

%%%%%%%%%%%%%%%%%%%%%%%%%%%%%%%%%%%%%%%%%%%%%%%%%%%%%%%%%%%%%%%%%%%%%%%%%%%%%%%%

The wall distance for each face point is used to interpolate the mean flow and
Reynolds stress variables from the profile data to the face point.
%
The Cholesky lower-triangular decomposition of the Reynolds stress tensor
is performed resulting in the following matrix being defined at each face point
%
\begin{equation*}
  L_{ij} = 
  \left[\begin{array}{ccc}
    \sqrt{R_{11}} & 0 & 0 \\[1.0ex]
    \frac{R_{21}}{L_{11}} & \sqrt{R_{22} - L^2_{21}} & 0 \\[1.0ex]
    \frac{R_{31}}{L_{11}} & \frac{R_{32} - L_{21} L_{31}}{L_{22}} & 
              \sqrt{R_{33} - L^2_{31} - L^2_{32}}
  \end{array}\right]
\end{equation*}

%###############################################################################

\subsubsubsection{Populate virtual box with randomly placed synthetic eddies}
\label{SEC:SEM:theory:AddEddies}

The SEM box is populated with randomly located synthetic eddies, with the
general idea shown in \fref{FIG:SEM:InflowBoxEddies}.
%
\begin{figure}[tb]
  \insertimage[width=0.95\textwidth]{sem/sem_eddies}
  \caption[Illustration of the SEM synthetic eddies.]%
  {Illustration of the SEM synthetic eddies randomly placed within the
  SEM virtual box. The red arrows represent the random intensity\slash spin
  for the eddy in the streamwise direction.}
  \label{FIG:SEM:InflowBoxEddies}
\end{figure}
%
Each synthetic eddy is not truly an eddy in the physical sense.
%
In reality, they are just source points that contain a randomized
\quoterm{intensity} in each coordinate direction.
%
Using a random number generator that produces a value $\sigma$ between 0 and 1,
each intensity is simply
%
\begin{equation*}
  \epsilon_j = 
  \begin{cases}
    -1 & \textrm{if } \sigma\leq 0.5\\
    +1 & \textrm{if } \sigma > 0.5
  \end{cases}
\end{equation*}
%
A random number is generated separately for each coordinate direction to get the
$\pm 1$ intensity value in each direction.
%
Since the term \quoterm{intensity} often implies the scale of a variable
instead of a binary value, it might be better to think of this as a
positive\slash negative \quoterm{spin} for the eddy in each direction.

%%%%%%%%%%%%%%%%%%%%%%%%%%%%%%%%%%%%%%%%%%%%%%%%%%%%%%%%%%%%%%%%%%%%%%%%%%%%%%%%

The input value for the SEM length scale, \lSEM, determines the
\quoterm{zone of influence} for each eddy.
%
This zone of influence is used in a distance-weighted shape function, described
in \sref{SEC:SEM:theory:SumFluctuations}, to determine how much an eddy affects
nearby face points.
%
Any face point that is beyond the distance \lSEM{} from the eddy will not be
affected by the eddy at all.

%%%%%%%%%%%%%%%%%%%%%%%%%%%%%%%%%%%%%%%%%%%%%%%%%%%%%%%%%%%%%%%%%%%%%%%%%%%%%%%%

The volume for an individual eddy,
\key[key:EddyVolume]{\vEddy}[eddy volume],
can be thought of as the eddy's total zone of influence.
%
By default, the zone of influence is applied separately for each coordinate
direction, so the volume of an eddy is approximated as a cube, $\vEddy=\lSEM^3$.
%
The input
\inp[nmlSEM::UseSphericalEddyVolume]{Use\_Spherical\_Eddy\_Volume}
can be used to treat the volume of the eddy as spherical,
$\vEddy = \sfrac{4}{3} \; \pi \left( \lSEM / 2 \right)^3
        = \sfrac{1}{6} \; \pi \: \lSEM^3$,
where \lSEM{} is the diameter of the sphere.

%%%%%%%%%%%%%%%%%%%%%%%%%%%%%%%%%%%%%%%%%%%%%%%%%%%%%%%%%%%%%%%%%%%%%%%%%%%%%%%%

Finally, the number of eddies, $N$, placed in an individual SEM box is
dependent on
\key[key:SEMVolume]{\vSEM}[Synthetic Eddy Method (SEM)!volume] and
\key[key:EddyVolume]{\vEddy}[eddy volume] using the equation
%
\begin{equation*}
  N = [\inp[nmlSEM::BoxNumEddiesMultiplier]{Box\_Num\_Eddies\_Multiplier}]
      \times \frac{\vSEM}{\vEddy}
\end{equation*}
%
The multiplier is an input that can be used to globally increase or decrease the
density of eddies within all of the SEM volumes.

%###############################################################################

\subsubsubsection{Sum SEM fluctuations at each face point}
\label{SEC:SEM:theory:SumFluctuations}

Velocity fluctuations are imposed at each face point on the SEM boundary by
summing the contributions from each synthetic eddy on that face point.
%
Let $u^*_{j,k}$ be the unscaled fluctuating velocity in direction $j$ for face
point $k$.
%
The contribution of the synthetic eddy $n$ towards $u^*_{j,k}$ is given by
\begin{equation*}
  u^*_{j,k,n} = \epsilon_j \tentfunc\left(\frac{x_k - x_n}{\lSEM}\right)
                           \tentfunc\left(\frac{y_k - y_n}{\lSEM}\right)
                           \tentfunc\left(\frac{z_k - z_n}{\lSEM}\right)
\end{equation*}
%
where
$\epsilon_j$ is the ``intensity'' or spin in the $j$-coordinate direction, and
$\tentfunc(a)$ is a distance-weighted shape function. 
%
GFR uses a tent function (hence the \tentsymbol) which is defined as
%
\begin{equation*}
  \tentfunc(a) = 
  \begin{cases}
    \sqrt{\tfrac{3}{2}}\left(1-\abs{a}\right) & \text{if } \abs{a} < 1 \\
    0                                         & \text{otherwise}
  \end{cases}
\end{equation*}
%
The total unscaled fluctuating velocity vector for face point $k$ is
obtained by summing the influence of all of the eddies on that point.
%
\begin{equation*}
  u^*_{j,k} = \sum_{n=1}^{N} u^*_{j,k,n}
\end{equation*}

%###############################################################################

\subsubsubsection{Compute BC states using SEM fluctuations}
\label{SEC:SEM:theory:ComputeBC}

Our ultimate goal is to compute the interior common fluxes from
\sref{SEC:Intro:FR:Upwind} for each face point on the SEM boundary.
%
Since we already know the interior state at each face point (found by evaluating
the solution polynomial from the interior cell at the face point), all that
remains is to determine the boundary\slash exterior state for these points.
%
Once the unscaled fluctuating velocity has been computed for face point $k$ as
shown in the previous section, it can be used with the SEM profile
data that was interpolated to this face point in
\sref{SEC:SEM:theory:DefineProfile}.
%
The scaled fluctuating velocity in coordinate direction $i$, $u'_i$, is computed
as the matrix multiplication
%
\begin{equation*}
  u'_{i,k} = L_{ij,k} u^*_{j,k}
\end{equation*}
%
Finally, the velocity for the boundary state, $u_{\subBC,i}$, in the coordinate
direction $i$ is the sum of the mean and fluctuating parts
%
\begin{equation*}
  u_{\subBC,i} = \overline{u}_{i} + u'_i
\end{equation*}
%
where each of these terms is assumed to be unique for the current face point $k$
so the subscript $k$ has been removed from each term for clarity.

%%%%%%%%%%%%%%%%%%%%%%%%%%%%%%%%%%%%%%%%%%%%%%%%%%%%%%%%%%%%%%%%%%%%%%%%%%%%%%%%

If the namelist input variable
\inp[nmlSEM::IncludeThermoFluctuations]{Include\_Thermo\_Fluctuations}%
=\texttt{.FALSE.} is used, then
SEM fluctuations will not be applied to the thermodynamic variables from
the SEM profile.
%
The remaining two boundary state variables,
%If SEM fluctuations are not being applied to the thermodynamic variables from
%the SEM profile, the remaining two boundary state variables,
$\rho_\subBC$ and $T_\subBC$, are set to the
value of the respective variable interpolated from the profile to the wall
distance of the face point, \thatis
%
\begin{align*}
  \rho_\subBC & = \overline{\rho}_k \\
  T_\subBC & = \overline{T}_k
\end{align*}
%
If \inp[nmlSEM::IncludeThermoFluctuations]{Include\_Thermo\_Fluctuations}%
=\texttt{.TRUE.}, then
SEM fluctuations will be applied to the thermodynamic variables.
%
These fluctuations are based on the Strong Reynolds Analogy following
equations 38 and 39 from Touber and Sandham\cite{Touber_2008}.
%
After some simplifications, the temperature fluctuations are given by
\begin{equation*}
  T' = \frac{- u' \overline{u}}{c_p}
\end{equation*}
%
Here, $u$ is the component of velocity in the streamwise direction and $c_p$ is
the specific heat at constant pressure.
%
The boundary state temperature is just the sum of the mean and fluctuating
temperatures
%
\begin{equation*}
  T_\subBC = \overline{T} + T'
\end{equation*}
%
and the boundary state density is then evaluated as
%
\begin{equation*}
  \rho_\subBC = \overline{\rho}\left(1-\frac{T'}{\overline{T}}\right)
\end{equation*}

%###############################################################################

\subsubsubsection{Convect eddies through the box}
\label{SEC:SEM:theory:ConvectEddies}

The last part of the SEM is convecting the eddies through the virtual box
(although, technically, this is the first step of the SEM for every time step
beyond the initialization).
%
All eddies are convected through the SEM box based on the SEM
convective velocity, \uSEM, specified in the input file.
%
With $x_n (t)$ giving the streamwise location of eddy $n$ at time $t$,
and $\Delta t$ giving the change in time between the current and
previous time steps, the new streamwise location of the eddy is
%
\begin{equation*}
  x_n (t+\Delta t) = x_n (t) + \uSEM \Delta t
\end{equation*}
%
If the new location for any eddy takes it outside the box, \thatis
$x_n>x_\subBC+\lSEM$, that respective eddy is deleted, and a new eddy is created
at a random location on the inflow of the SEM box to take the place of the old
eddy.
%
Because it is highly probable that multiple eddies will leave the box each time
step, a small randomized offset is added to the streamwise location of each new
eddy to prevent \quoterm{waves of eddies} that share the same axial location in
the SEM box.

%###############################################################################

